\section{A single detecting correlation and an exact occurrence reduction}\label{sec:probe}
\subsection{A compact test which misses no nontrivial zero}
Let $a_j=2^{-j}$ and let b be the density of the absolutely
convergent sum of independent uniforms on $[-a_j,a_j]$.
Its bilateral Laplace transform is
\begin{equation}\label{eq:probe-product}
\mathcal B(z)=\prod_{j\ge1}\frac{\sinh(2^{-j}z)}{2^{-j}z}.
\end{equation}
Removable factors at zero have value one. Define
\begin{equation}\label{eq:probe}
h_*=\frac{e^x b(x)}{\mathcal B(1)},\qquad f_*=\T h_*,
\qquad
F(z)=\int e^{zx}f_*(x)\dd x
=(1/4-z^2)\frac{\mathcal B(1+z)}{\mathcal B(1)} .
\end{equation}

\begin{lemma}\label{lem:probe}
The real functions $h_*,f_*$ are smooth, supported in
$[-1,1]$, and $f_*\in\cD^0$. The entire function F has
no zeros in $|\Ree z|<1/2$. In particular
$\widehat f_*(\tau)\ne0$ at every real $\tau$.
\end{lemma}
\begin{proof}
The random sum is supported in $[-1,1]$. On the imaginary
axis, bound any fixed number k of product factors by
$\min(1,(a_j|\tau|)^{-1})$ and the rest by one.
Its Fourier transform decays faster than every power,
giving a smooth compact density by Fourier inversion.
The product converges locally uniformly in the plane
because its tail deviations from one are summable.
All zeros of each factor lie on the imaginary axis, and
off that axis the nonzero summable tail cannot vanish.
For $|\Ree z|<1/2$, one has $\Ree(1+z)>1/2$, so neither
the product nor $1/4-z^2$ vanishes. Integration by parts
proves \eqref{eq:probe} and the two neutral moments.
\end{proof}
The tilt $e^x$ removes the real Fourier zeros that an
untilted even compact density would have. This probe
uses no zero list and no assumption about RH.

Put $k(t)=\ip{f_*}{U_tf_*}_2$ and define the real even response
\begin{align}\label{eq:C-star}
C_*(t)=Q(f_*,U_tf_*)&=
\frac1{2\pi}\int m_+(\tau)|\widehat f_*(\tau)|^2
                                     e^{-\ii t\tau}\dd\tau\\
&\quad-\sum_{a\ge2}\frac{\Lambda(a)}{\sqrt a}
                  \{k(t+\log a)+k(t-\log a)\}.\nonumber
\end{align}
Here k is smooth and supported in $[-2,2]$, so the prime
sum is locally finite in t. The definition contains no
spectral zero data. For $t>2$ its archimedean part is
\begin{equation}
C_{*,\infty}(t)=
-\sum_{j\ge1}e^{-(2j+1/2)t}
                   F(2j+1/2)F(-2j-1/2).
\end{equation}
This follows by expanding the off-diagonal kernel.
The $j=0$ term vanishes by pole neutrality, and the local
contact contributes only for overlapping supports.

\subsection{An unconditional criterion and an unconditional nondecay theorem}
To avoid confusion with the plaquette density $\rho$, denote
a nontrivial zeta zero by $\varrho$, counted with multiplicity,
and put $\lambda_\varrho=\varrho-1/2$.
The unconditioned explicit formula \cite{cc} gives
\begin{equation}\label{eq:C-zeros}
C_*(t)=\sum_\varrho c(\lambda_\varrho)e^{\lambda_\varrho t},
\qquad c(\lambda)=F(\lambda)F(-\lambda).
\end{equation}
Opposite Fourier conventions are reconciled by
$\varrho\mapsto1-\varrho$. Off the critical line the
coefficient is this product, not a squared modulus.
It is nonzero at every nontrivial zero by Lemma~\ref{lem:probe}.
The zero symmetries, $0<\Ree\varrho<1$, and
$N(T)=O(T\log T)$ are unconditional \cite{dlmf,elkies}.
Smooth compact support gives rapid vertical decay of F
uniformly in bounded real strips. Thus the sum converges
absolutely and locally uniformly with every t derivative.

\begin{theorem}[One-probe regularity criterion]\label{thm:probe-criterion}
The following are equivalent: RH; boundedness of $C_*$ on
$\R$; temperedness of the distribution defined by $C_*$;
and, for every $\epsilon>0$, a bound
$|C_*(t)|\le C_\epsilon e^{\epsilon t}$ for $t\ge0$.
Moreover
\begin{equation}\label{eq:growth-rate}
\limsup_{t\to+\infty}\frac{\log(1+|C_*(t)|)}{t}
=\sup_\varrho|\Ree\varrho-1/2|.
\end{equation}
\end{theorem}
\begin{proof}
Under RH the expansion becomes
\begin{equation}\label{eq:C-positive}
C_*(t)=\sum_{\gamma\ {\rm distinct}}
m_\gamma|\widehat f_*(\gamma)|^2e^{-\ii\gamma t}.
\end{equation}
The positive coefficients are summable; this proves all
three regularity properties.

Conversely, in $\Ree z>1/2$ absolute convergence permits
the one-sided Laplace transform
\begin{equation}\label{eq:laplace-probe}
\int_0^\infty e^{-zt}C_*(t)\dd t
=\sum_\varrho\frac{c(\lambda_\varrho)}{z-\lambda_\varrho}.
\end{equation}
The right side is meromorphic on the plane, normally
convergent off its discrete poles; the residue at a zero
of multiplicity m is $m c(\lambda)\ne0$.
A subexponential bound makes the left side holomorphic
on $\Ree z>0$. Temperedness has the same consequence:
multiply the smooth function distribution by a smooth
cutoff equal to one for $t\ge1$ and zero for $t\le0$,
take its exponentially damped distributional Laplace
transform, and add the compact initial-interval correction.
This avoids multiplying an arbitrary distribution by a
discontinuous step. The transform agrees with the integral
in its initial half-plane and is holomorphic for $\Ree z>0$.
It cannot have one of the nonzero residues there.
The functional equation excludes left-half-plane
$\lambda_\varrho$ as well. Boundedness is a special case.

For \eqref{eq:growth-rate} let b be the supremum on the
right, also the supremum of $\Ree\lambda_\varrho$.
Absolute summability bounds $|C_*(t)|$ by
$e^{bt}\sum|c(\lambda_\varrho)|$. If the limsup were
less than b, a slightly larger exponent still below b
would give a holomorphic Laplace transform to the left
of at least one nonzero-residue pole, a contradiction.
For b=0 the upper bound and the nonnegative logarithm
give equality directly.
\end{proof}
This is a relation between RH and a precise regularity
property, not a derivation of that property from OS
temperedness. The variable t is logarithmic character
scale, and the response has not been shown to be a
Schwinger function in spacetime. Simply postulating its
temperedness would postulate an RH-equivalent arithmetic
condition without establishing a physical explanation.

\begin{corollary}\label{cor:nondecay}
Unconditionally, $C_*\notin C_0(\R)$.
\end{corollary}
\begin{proof}
Membership in $C_0$ would imply boundedness and hence RH.
Under that consequence, the absolutely convergent series
\eqref{eq:C-positive} is uniformly almost periodic and
\[
\lim_{T\to\infty}\frac1{2T}\int_{-T}^T|C_*(t)|^2\dd t
=\sum_{\gamma\ {\rm distinct}}
\{m_\gamma|\widehat f_*(\gamma)|^2\}^2>0.
\]
One proves the formula first for finite sums and passes
by uniform convergence. There are nontrivial zeros and
the probe misses none. A function in $C_0$ has zero such
mean square, a contradiction.
\end{proof}
Together with Theorem~\ref{thm:current-spectrum}, this
excludes $C_*(t)=\ip\xi{\UU_t\eta}$ for \emph{any} two
physical Hilbert vectors. In particular no current-covariant
source law with the exact Weil pairing can exist in the
fixed boundary space, even if its seed is not a class
function or a finite-order generalized source.
An inserted coefficient $\ip v{K\UU_tv}$ is also excluded
whenever $v\in\Dom K^*$, because it is the ordinary
coefficient $\ip{K^*v}{\UU_tv}$.
A genuinely distributional left functional is outside
this finite-vector statement and is not an OS norm.

Similarly, a diagonal cutoff $A(R)$ cannot produce
bounded actual response vectors with all translated
mixed limits equal to $C_*$. A weakly convergent
subsequence of such bounded vectors, paired with the
strongly recovered translates in \eqref{eq:rec-pair},
would give a finite-vector current coefficient equal
to $C_*$, contradicting the corollary. This is a
conditional obstruction to that proposed joint limit,
not a claim that the joint convergence holds.

\subsection{One actual stationary correlation would suffice}
\begin{theorem}[One-correlation source occurrence]\label{thm:one-occurrence}
Suppose an independently specified physical closed
admissible sector contains a vector v and a strongly
continuous unitary group $V_t=e^{-\ii tG}$ such that
\begin{equation}\label{eq:one-occurrence-hyp}
\ip v{V_tv}_{\rm phys}=C_*(t)\quad(t\in\R).
\end{equation}
Then RH follows, and spectral filtering in that same
sector constructs a continuous global source
$J:\cD^0\to\cH_{\rm phys}$ with
$Jf_*=v$, $JU_t=V_tJ$, and
$\ip{Jf}{Jg}_{\rm phys}=Q(f,g)$.
\end{theorem}
\begin{proof}
The left side of \eqref{eq:one-occurrence-hyp} is bounded
by $\norm v^2$, so Theorem~\ref{thm:probe-criterion} gives
RH. Fourier uniqueness for finite measures and
\eqref{eq:C-positive} determine the actual spectral
measure of v:
\begin{equation}
\dd\mu_v(\lambda)=
\sum_{\gamma\ {\rm distinct}}
m_\gamma|\widehat f_*(\gamma)|^2\delta_\gamma(\dd\lambda).
\end{equation}
For $g\in\cD^0$ define
\begin{equation}\label{eq:filter-source}
Jg=\left(\frac{\widehat g}{\widehat f_*}\right)(G)v .
\end{equation}
There are no real denominator zeros, but no polynomial
bound for this quotient is presumed. Its actual
functional-calculus domain is verified instead:
\[
\int\left|\frac{\widehat g(\lambda)}
                    {\widehat f_*(\lambda)}\right|^2\dd\mu_v
=\sum_\gamma m_\gamma|\widehat g(\gamma)|^2<\infty.
\]
Thus bounded spectral truncations converge to an actual
vector in the same closed sector. Polarization gives
\eqref{eq:zero-pairing}, hence Q by the explicit formula.
The ratio for $g=f_*$ is one; multiplication by
$e^{-\ii t\lambda}$ proves covariance. Finally for tests
in a fixed compact interval, integration by parts bounds
$|\widehat g(\gamma)|$ by
$C_{I,M}(1+|\gamma|)^{-M}$ times a fixed finite derivative
seminorm. Zero counting makes the squared sum finite,
proving test-space continuity.
\end{proof}
An equivalent hypothesis is an actual norm-continuous
stationary family $\{v_t\}$ with
$\ip{v_s}{v_t}=C_*(t-s)$: translation on its linear
span is an isometry, extends to a unitary group on its
closure, and gives \eqref{eq:one-occurrence-hyp}.
A separately identified native arithmetic generator is
therefore optional. No family with this correlation has
been constructed here. The theorem guarantees vectors
in the declared closed sector, not bounded local
operators or stronger electric source bounds unless
those additional domains are independently preserved.
